% GENERATED FILE. Edit canonical lessons, then run npm run generate:books.
% canonical-source-hash: fnv1a64:a8e4c73a565d97e0
% canonical-lessons: JA-C302-yuurei, JA-C302-obake, JA-C302-uramu, JA-C302-densetsu, JA-C302-kamisama

\chapter{Ghost Stories}
\label{ch:ja-ghost-stories}
\hlchaptermodality{302}

\begin{chapteropening}
\textbf{By the end of this chapter:} \emph{I can tell what the ghost in the old story wanted.}

Complete the last lesson of chapter 302: I can tell what the ghost in the old story wanted.
\end{chapteropening}

\section[\texorpdfstring{\ja{ゆうれい} (yūrei)}{ゆうれい (yūrei)}]{\ja{ゆうれい} (yūrei) — a ghost}

\label{lesson:JA-C302-yuurei}



\begin{quote}
Before the new one: say the Japanese for poor, then the Japanese for greedy.
\end{quote}

\subsection*{You'll want to know: \ja{ゆうれい}}
\textbf{\ja{ゆうれい}} — \emph{yūrei} — \textquotedblleft{}a ghost\textquotedblright{}.

\begin{quote}
\textbf{\ja{その} \ja{ふるい} \ja{てらには} \ja{ゆうれいが} \ja{います。}} — \emph{sono furui tera ni wa yūrei ga imasu} — \textquotedblleft{}There is a ghost in that old temple.\textquotedblright{}
\end{quote}

\begin{cousinweb}
\textbf{\ja{ゆうれい}} (\emph{yūrei}) is Sino-Japanese (kango): two Chinese characters read with their on'yomi, readings borrowed from Chinese. They are \emph{yū}, \textquotedblleft{}faint\textquotedblright{} (Middle Chinese \emph{'jiw}) and \emph{rei}, \textquotedblleft{}spirits\textquotedblright{} (\emph{leng}) (KANJIDIC2; Middle Chinese in Baxter's transcription).
\end{cousinweb}

\begin{grammarlens}[title={What you've built}]
One more word for telling what happened.
\end{grammarlens}

\subsection*{Your turn}
\begin{itemize}
\raggedright
  \item \emph{Say it:} \emph{yūrei}
  \item \emph{Say it:} \emph{yūrei}, once more
  \item \emph{Say it:} \emph{sono furui tera ni wa yūrei ga imasu} — There is a ghost in that old temple.
  \item \emph{Recall:} say the Japanese for poor, then the Japanese for greedy, then say \emph{yūrei} again
\end{itemize}

\begin{tcolorbox}[breakable,colback=teal!4,colframe=teal!35!black,title={Before you move on}]
What does \ja{ゆうれい} mean? (\textquotedblleft{}A ghost\textquotedblright{}.) Say it once more.
\end{tcolorbox}
\section[\texorpdfstring{\ja{おばけ} (obake)}{おばけ (obake)}]{\ja{おばけ} (obake) — a spook; a monster}

\label{lesson:JA-C302-obake}



\begin{quote}
Before the new one: say the Japanese for greedy, then the Japanese for a ghost.
\end{quote}

\subsection*{You'll want to know: \ja{おばけ}}
\textbf{\ja{おばけ}} — \emph{obake} — \textquotedblleft{}a spook; a monster\textquotedblright{}.

\begin{quote}
\textbf{\ja{いもうとは} \ja{おばけが} \ja{こわいです。}} — \emph{imōto wa obake ga kowai desu} — \textquotedblleft{}My little sister is afraid of ghosts.\textquotedblright{}
\end{quote}

\begin{cousinweb}
\textbf{\ja{おばけ}} (\emph{obake}) is a native Japanese word (wago), inherited rather than borrowed from Chinese, built from \emph{o-}, \textquotedblleft{}the polite prefix\textquotedblright{}, and \emph{bake}, \textquotedblleft{}changing shape\textquotedblright{} (from \emph{bakeru}) (JMdict, KANJIDIC2).
\end{cousinweb}

\begin{grammarlens}[title={What you've built}]
One more word for telling what happened.
\end{grammarlens}

\subsection*{Your turn}
\begin{itemize}
\raggedright
  \item \emph{Say it:} \emph{obake}
  \item \emph{Say it:} \emph{obake}, once more
  \item \emph{Say it:} \emph{imōto wa obake ga kowai desu} — My little sister is afraid of ghosts.
  \item \emph{Recall:} say the Japanese for greedy, then the Japanese for a ghost, then say \emph{obake} again
\end{itemize}

\begin{tcolorbox}[breakable,colback=teal!4,colframe=teal!35!black,title={Before you move on}]
What does \ja{おばけ} mean? (\textquotedblleft{}A spook; a monster\textquotedblright{}.) Say it once more.
\end{tcolorbox}
\section[\texorpdfstring{\ja{うらむ} (uramu)}{うらむ (uramu)}]{\ja{うらむ} (uramu) — to bear a grudge against}

\label{lesson:JA-C302-uramu}



\begin{quote}
Before the new one: say the Japanese for a ghost, then the Japanese for a spook.
\end{quote}

\subsection*{You'll want to know: \ja{うらむ}}
\textbf{\ja{うらむ}} — \emph{uramu} — \textquotedblleft{}to bear a grudge against\textquotedblright{}.

\begin{quote}
\textbf{\ja{ゆうれいは} \ja{その} \ja{おとこを} \ja{うらんでいました。}} — \emph{yūrei wa sono otoko o urande imashita} — \textquotedblleft{}The ghost held a grudge against the man.\textquotedblright{}
\end{quote}

Polite \textbf{\ja{うらみます}}, \emph{uramimasu}; a godan verb. Past (polite): \textbf{\ja{おんなが} \ja{おとこを} \ja{うらみました。}}, \emph{onna ga otoko o uramimashita}.

\begin{cousinweb}
\textbf{\ja{うらむ}} (\emph{uramu}) is a native Japanese word (wago), inherited rather than borrowed from Chinese; its kanji was attached later as a kun'yomi (JMdict).
\end{cousinweb}

\begin{grammarlens}[title={What you've built}]
One more word for telling what happened.
\end{grammarlens}

\subsection*{Your turn}
\begin{itemize}
\raggedright
  \item \emph{Say it:} \emph{uramu}
  \item \emph{Say it:} \emph{uramu}, once more
  \item \emph{Say it:} \emph{yūrei wa sono otoko o urande imashita} — The ghost held a grudge against the man.
  \item \emph{Recall:} say the Japanese for a ghost, then the Japanese for a spook, then say \emph{uramu} again
\end{itemize}

\begin{tcolorbox}[breakable,colback=teal!4,colframe=teal!35!black,title={Before you move on}]
What does \ja{うらむ} mean? (\textquotedblleft{}To bear a grudge against\textquotedblright{}.) Say it once more.
\end{tcolorbox}
\section[\texorpdfstring{\ja{でんせつ} (densetsu)}{でんせつ (densetsu)}]{\ja{でんせつ} (densetsu) — a legend}

\label{lesson:JA-C302-densetsu}



\begin{quote}
Before the new one: say the Japanese for a spook, then the Japanese for to bear a grudge against.
\end{quote}

\subsection*{You'll want to know: \ja{でんせつ}}
\textbf{\ja{でんせつ}} — \emph{densetsu} — \textquotedblleft{}a legend\textquotedblright{}.

\begin{quote}
\textbf{\ja{この} \ja{やまには} \ja{ふるい} \ja{でんせつが} \ja{あります。}} — \emph{kono yama ni wa furui densetsu ga arimasu} — \textquotedblleft{}There is an old legend about this mountain.\textquotedblright{}
\end{quote}

\begin{cousinweb}
\textbf{\ja{でんせつ}} (\emph{densetsu}) is Sino-Japanese (kango): two Chinese characters read with their on'yomi, readings borrowed from Chinese. They are \emph{den}, \textquotedblleft{}transmit\textquotedblright{} (Middle Chinese \emph{drjwen}) and \emph{setsu}, \textquotedblleft{}explanation\textquotedblright{} (\emph{sywet}) (KANJIDIC2; Middle Chinese in Baxter's transcription).
\end{cousinweb}

\begin{grammarlens}[title={What you've built}]
One more word for telling what happened.
\end{grammarlens}

\subsection*{Your turn}
\begin{itemize}
\raggedright
  \item \emph{Say it:} \emph{densetsu}
  \item \emph{Say it:} \emph{densetsu}, once more
  \item \emph{Say it:} \emph{kono yama ni wa furui densetsu ga arimasu} — There is an old legend about this mountain.
  \item \emph{Recall:} say the Japanese for a spook, then the Japanese for to bear a grudge against, then say \emph{densetsu} again
\end{itemize}

\begin{tcolorbox}[breakable,colback=teal!4,colframe=teal!35!black,title={Before you move on}]
What does \ja{でんせつ} mean? (\textquotedblleft{}A legend\textquotedblright{}.) Say it once more.
\end{tcolorbox}
\section[\texorpdfstring{\ja{かみさま} (kamisama)}{かみさま (kamisama)}]{\ja{かみさま} (kamisama) — a god; a deity}

\label{lesson:JA-C302-kamisama}



\begin{quote}
Before the new one: say the Japanese for to bear a grudge against, then the Japanese for a legend.
\end{quote}

\subsection*{You'll want to know: \ja{かみさま}}
\textbf{\ja{かみさま}} — \emph{kamisama} — \textquotedblleft{}a god; a deity\textquotedblright{}.

\begin{quote}
\textbf{\ja{そふは} \ja{やまの} \ja{かみさまに} \ja{いのりました。}} — \emph{sofu wa yama no kamisama ni inorimashita} — \textquotedblleft{}Grandfather prayed to the god of the mountain.\textquotedblright{}
\end{quote}

\begin{cousinweb}
\textbf{\ja{かみさま}} (\emph{kamisama}) is a native Japanese word (wago), inherited rather than borrowed from Chinese, built from \emph{kami}, \textquotedblleft{}god\textquotedblright{}, and \emph{sama}, \textquotedblleft{}the honorific\textquotedblright{} (JMdict, KANJIDIC2).
\end{cousinweb}

\begin{grammarlens}[title={What you've built}]
One more word for telling what happened.
\end{grammarlens}

\subsection*{Your turn}
\begin{itemize}
\raggedright
  \item \emph{Say it:} \emph{kamisama}
  \item \emph{Say it:} \emph{kamisama}, once more
  \item \emph{Say it:} \emph{sofu wa yama no kamisama ni inorimashita} — Grandfather prayed to the god of the mountain.
  \item \emph{Recall:} say the Japanese for to bear a grudge against, then the Japanese for a legend, then say \emph{kamisama} again
\end{itemize}

\begin{tcolorbox}[breakable,colback=teal!4,colframe=teal!35!black,title={Before you move on}]
What does \ja{かみさま} mean? (\textquotedblleft{}A god; a deity\textquotedblright{}.) Say it once more.
\end{tcolorbox}
